\section{Evidence Synthesis}
\label{sec:evidence-synthesis}

This section synthesizes the Core Evidence Corpus defined in Section~\ref{sec:methodology}, using contextual references only to interpret mechanisms and substrates. The main pattern is clear: direct VLA attack evidence is now substantial enough to review, but it remains concentrated in OpenVLA/LIBERO-style manipulation, often preprint-heavy, and uneven in physical validation. Embodied-adjacent attacks remain useful for threat modeling, while contextual references should not be written as if they were validated VLA attacks.

\begin{table*}[t]
\centering
\scriptsize
\caption{Evidence-aware synthesis by corpus. Core rows are attack evidence; contextual rows support taxonomy or evaluation but are not included in attack-study denominators. D: digital; S: simulation; HIL: hardware-in-the-loop; P: physical robot/world; PR: peer-reviewed; PP: preprint or non-archival version.}
\label{tab:attack-literature-summary}
\renewcommand{\arraystretch}{1.1}
\setlength{\tabcolsep}{1.8pt}
\begin{tabularx}{\textwidth}{@{}p{0.13\textwidth}p{0.28\textwidth}p{0.10\textwidth}p{0.11\textwidth}Y@{}}
\toprule
\textbf{Corpus / stratum} & \textbf{Representative work} & \textbf{Status} & \textbf{Validation} & \textbf{What the evidence supports} \\
\midrule
Direct VLA: instruction/control & RoboGCG, VLA-Fool, SABER, Q-DIG, trajectory redirection \cite{jones2025robogcg,yan2025vlafool,wu2026saber,srikanth2026qdig,puthumanaillam2026trajectory} & PP & D/S/P & Language or multimodal perturbations can alter VLA rollout behavior; black-box, natural-prompt, and command-preserving settings are emerging. \\
Direct VLA: patches/physical vision & ICCV patch attacks, EDPA, ADVLA, UPA-RFAS, VLA-Hijack, Tex3D, partially observable patches \cite{wang2025vlaadversarial,xu2025edpa,zhang2025advla,lu2026patch,fu2026vlahijack,chen2026tex3d,wang2026partialpatch} & PR/PP & D/S/P & Visual attacks can target embeddings, grounding, visual proprioception, action trajectories, and object textures; transfer and physical evidence are improving but uneven. \\
Direct VLA: action/backdoor & FreezeVLA, BadVLA, GoBA, DropVLA/TabVLA, AttackVLA/BackdoorVLA, SilentDrift \cite{wang2025freezevla,zhou2025badvla,zhou2025goba,xu2025dropvla,li2025attackvla,xu2026silentdrift} & PP & S/P & VLA attacks can target action primitives, action freezing, chunk drift, and long-horizon triggered behavior, often while preserving clean utility. \\
Direct VLA: reasoning/privacy & TRAP, altered reasoning traces, ReasonBreak, VLA membership inference, VLALeaks \cite{huang2026trap,altered2026thoughts,teymoorianfard2026reasonbreak,peng2026miavla,luan2026vlaleaks} & PP & D/S/P & Reasoning traces and training membership are VLA-specific target representations; objectives include integrity and confidentiality rather than only task failure. \\
Core: embodied-adjacent attacks & BadRobot, RoboPAIR, CHAI, RoboJailBench, ANNIE, Phantom Menace, Eva-VLA \cite{zhang2025badrobot,robey2025robopair,burbano2025chai,yeke2026robojailbench,huang2025annie,lu2025phantom,liu2025evavla} & PR/PP & D/S/P & Embodied LLM/VLM, environmental command, trajectory, sensor, and physical-variation attacks inform VLA threat models but are not counted as direct VLA evidence unless VLA policies are attacked. \\
Contextual references & VLA/RFM substrates, benchmarks, adversarial ML, physical/sensor attacks, RAG/tool-agent security, RL backdoors, ROS/robot cybersecurity \cite{kim24openvla,black2024pi0,liu2023libero,goodfellow2015explaining,brown2017adversarial,cao2019lidar,zou2025poisonedrag,debenedetti2024agentdojo,kiourti2020trojdrl,deng2022insecurityros2} & PR/PP & D/S/P & Provides mechanisms, substrates, and evaluation lessons; not included in attack-study denominators unless the work has explicit attacker, target, method, and outcome. \\
\bottomrule
\end{tabularx}
\end{table*}


\subsection{Training-Time and Supply-Chain Attacks}

\textbf{Direct evidence.} VLA backdoor work has expanded from generic task disruption to action and trajectory control. BadVLA studies objective-decoupled VLA backdoors that preserve clean performance while inducing triggered action deviations \cite{zhou2025badvla}. Its evidence is direct because the victim is a VLA policy and the measured output is action behavior, but its strongest claims are still tied to the particular trigger, benchmark, and fine-tuning setup. GoBA changes the threat model by using ordinary physical objects as goal-oriented triggers and by introducing BadLIBERO-style evaluation with ``nothing to do,'' ``try to do,'' and ``success to do'' levels \cite{zhou2025goba}. This is stronger than a binary failure label because it distinguishes a triggered attempt from completed attacker-specified behavior.

DropVLA/TabVLA narrows the target further to reusable low-level action primitives, such as gripper release, under small poisoning budgets and reports simulation plus physical feasibility \cite{xu2025dropvla}. Its main contribution for a robotics audience is temporal precision: the attack is not only that a task fails, but that a primitive fires at a decision window. AttackVLA then provides a broader evaluation harness for adversarial and backdoor attacks and introduces BackdoorVLA for attacker-specified long-horizon action sequences \cite{li2025attackvla}. SilentDrift is especially robotics-relevant because it attacks action chunking and delta-pose integration, showing how small poisoned drifts can accumulate during intra-chunk open-loop execution \cite{xu2026silentdrift}. Together, these papers show a progression from \emph{triggered failure} to \emph{triggered goal behavior}, \emph{primitive-level control}, and \emph{chunk-level trajectory drift}. They should not be collapsed into one backdoor category.

\textbf{Evidence strength and gaps.} The direct VLA backdoor evidence is stronger than analogies from classifiers because it measures action traces and clean-task retention. It is still not yet field-wide evidence. Most studies use curated manipulation suites and controlled triggers; few report how data inspection, trajectory validation, controller clipping, or post-training updates affect attack survival. Few compare whether the same trigger works across autoregressive token heads, continuous action heads, action chunks, and flow/diffusion policies. These gaps matter because training-time attacks become deployment risks only if the trigger persists through the actual adaptation pipeline, robot morphology, camera placement, and controller.

\textbf{Embodied-adjacent and analogical evidence.} Classic poisoning and backdoor attacks establish the general supply-chain risk \cite{biggio2012poisoning,gu2017badnets,liu2018trojaning,carlini2024poisoningweb}, and RL backdoors show that sequential policies can learn trigger-conditioned behavior \cite{kiourti2020trojdrl,wang2021backdoorl,cui2024badrl,rathbun2025inception}. These works justify VLA backdoor threat models but do not by themselves prove VLA vulnerability. The direct VLA backdoor papers now provide that evidence, with the caveat that most remain recent preprints and frequently rely on a small set of benchmarks and robot embodiments.

\subsection{Perception, Multimodal, and Physical Attacks}

\textbf{Direct evidence.} Early VLA patch work formulates action-discrepancy, position-aware, and targeted manipulation objectives and evaluates digital plus physical patch settings \cite{wang2025vlaadversarial}. This work is important because the objective is defined in action space rather than as image classification error. VLA-Fool expands the attack surface to textual, visual, and cross-modal misalignment attacks under white-box and black-box settings \cite{yan2025vlafool}. Its strongest value is not a single failure number, but the comparison across language perturbations, visual patches/noise, and image--instruction mismatch.

EDPA reduces the required attacker knowledge by targeting visual encoder embeddings and cross-modal alignment rather than the full action head \cite{xu2025edpa}. ADVLA pushes a related gray-box direction by using attention-guided sparse feature-space perturbations that are low-amplitude and patch-wise localized \cite{zhang2025advla}. Universal transferable patch work, such as UPA-RFAS, studies black-box transfer across VLA models, prompts, viewpoints, and sim-to-real settings \cite{lu2026patch}. VLA-Hijack targets visual proprioception by suppressing real-arm features and injecting a phantom embodiment patch, exposing a representation-level reason for cross-architecture transfer \cite{fu2026vlahijack}. Tex3D moves from 2D patches to adversarial object textures optimized over long-horizon trajectories and evaluated in simulation and real-robot settings \cite{chen2026tex3d}. Partially observable patch attacks further constrain the adversary to a short trajectory prefix, a more realistic physical-access setting \cite{wang2026partialpatch}.

These papers form a useful sequence of evidence. Action-space patches ask whether visual perturbations can change robot motion. Embedding and attention attacks ask whether attacking the shared VLM/VLA representation is enough when the action head is unknown. Transferable patches ask whether the attack survives architecture, fine-tuning, and viewpoint changes. 3D texture attacks ask whether the attack can be bound to manipulated objects rather than pasted into the camera frame. Partial-observation attacks ask whether the adversary can attack without seeing the full future rollout. The sequence is more informative than a single ``visual attack'' bucket.

\textbf{Analogical evidence.} Foundational adversarial examples, physical patches, object-detector attacks, 3D mesh attacks, LiDAR attacks, and audio-command attacks show that physical sensor inputs can be manipulated \cite{szegedy2014intriguing,goodfellow2015explaining,papernot2016limitations,carlini2017towards,kurakin2017physical,brown2017adversarial,moosavi2017universal,madry2018towards,athalye2018synthesizing,eykholt2018robust,eykholt2018detectors,chen2018shapeshifter,thys2019fooling,xu2020tshirt,xiao2019meshadv,cao2019lidar,tu2020lidar,sun2020robustlidar,cao2021invisible,carlini2016hiddenvoice,zhang2017dolphinattack,yuan2018commandersong,sugawara2020lightcommands}. Their direct implication for VLA systems depends on whether the perturbation changes grounding, plan, action, trajectory, or recovery. This is why VLA patch papers should report action traces and consequence labels rather than only image-level success.

\subsection{Instruction, Planning, and Reasoning Attacks}

\textbf{Direct evidence.} RoboGCG adapts adversarial prompt optimization to VLA control authority and shows that a prompt-level intervention at rollout start can influence action reachability over time \cite{jones2025robogcg}. Its claim is stronger than a text jailbreak claim because the target is robot action reachability, but the access and optimization assumptions still need to be reported whenever the result is cited. SABER makes the instruction threat model more practical by using a black-box, agentic edit process with bounded character/token/prompt edits and rollout-level feedback across multiple VLA models \cite{wu2026saber}. Q-DIG approaches red-teaming differently: rather than optimizing a single adversarial suffix, it uses quality-diversity search to generate natural, task-relevant instructions that expose robot-policy failures and can be used for robustness improvement \cite{srikanth2026qdig}. Trajectory-level redirection studies command-preserving prompts that keep the instruction close to the intended task while redirecting the physical outcome \cite{puthumanaillam2026trajectory}.

Reasoning-enabled VLAs create an additional internal target. TRAP shows that CoT or intermediate reasoning traces can be hijacked by physical patches to steer downstream actions without changing the user's instruction \cite{huang2026trap}. Altered-thought studies further indicate that entity references in VLA reasoning traces can causally affect robot manipulation success, while non-reasoning control models may be immune to that specific channel \cite{altered2026thoughts}. ReasonBreak extends the reasoning-trajectory question to autonomous-driving-style VLA systems and evaluates black-box textual corruptions against reasoning and trajectory safety metrics in closed-loop simulation \cite{teymoorianfard2026reasonbreak}. The shared lesson is that reasoning text is neither merely explanation nor merely prompt context; in some VLA designs it is a control-relevant representation.

\textbf{Embodied-adjacent evidence.} BadRobot, RoboPAIR, CHAI, SafeAgentBench, EARBench, and RoboJailBench evaluate embodied jailbreak, physical command hijacking, and risk-aware planning where success is measured closer to action than text-only refusal \cite{zhang2025badrobot,robey2025robopair,burbano2025chai,yin2024safeagentbench,zhu2024earbench,yeke2026robojailbench}. These are important for goal hijacking and HRI threat models, but they should be labeled as embodied-adjacent when the target is not a VLA policy.

\subsection{Action-Level and Closed-Loop Attacks}

Action representation is the robotics core of this review. Autoregressive action-token policies expose token probabilities and discretized gripper/pose bins; a prompt or patch can change which token is selected, and early token errors can propagate through subsequent actions. RT-1, RT-2, OpenVLA, and SpatialVLA illustrate tokenized or discretized action interfaces whose action vocabulary, spatial grids, or detokenizers become attack-relevant representations \cite{brohan2022rt1,brohan2023rt2,kim24openvla,qu2025spatialvla}. Parallel action chunks create a different risk: the model commits to several future actions before new observations can correct the error. ACT, OpenVLA-OFT, and related high-throughput VLA fine-tuning recipes show why chunk length, temporal ensemble, and parallel decoding should be reported as security variables rather than only efficiency variables \cite{zhao2023act,kim2025openvlaoft}. SilentDrift directly exploits this intra-chunk open-loop property \cite{xu2026silentdrift}.

Waypoint and trajectory policies expose geometric objectives such as end-effector displacement, curvature, unsafe workspace entry, or targeted path redirection \cite{wang2025vlaadversarial,puthumanaillam2026trajectory}. Planner-mediated systems such as CLIPort, Code as Policies, ProgPrompt, and VoxPoser make object binding, generated code, value maps, and waypoint fields explicit enough to inspect but also explicit enough to attack \cite{shridhar2022cliport,liang2023code,singh2023progprompt,huang2023voxposer}. Diffusion and flow policies, including Diffusion Policy and $\pi_0$, shift the target from a token to a conditional trajectory distribution, denoising path, velocity field, or sampling process \cite{chi2023diffusion,black2024pi0}. Generalist policy and dataset work such as Octo, Open X-Embodiment, DROID, BridgeData, RoboNet, and Gato further matter because they define the cross-embodiment and data-provenance assumptions under which attacks may or may not transfer \cite{ghosh2024octo,openx2023,khazatsky2024droid,walke2023bridgedata,dasari2019robonet,reed2022gato}.

The closed-loop consequence depends on controller and recovery design. A temporal ensemble can damp high-frequency perturbations but may also prolong a poisoned action chunk. A high re-planning frequency can recover from one-frame perception attacks but may be vulnerable to repeated low-amplitude nudges. Controller latency and stale observations can turn a semantically correct command into an unsafe physical action. Safety envelopes and low-level controllers can prevent direct collision while still allowing wrong-object manipulation, unauthorized skill use, privacy leakage, or near-boundary behavior. Therefore, direct VLA attack papers should report not only ASR but also first affected representation, action deviation, recovery opportunity, controller intervention, and whether the safety envelope changed the outcome.

\subsection{System, Memory, Tool, HRI, and Persistent Attacks}

\textbf{Embodied-adjacent evidence.} Tool-agent attacks and indirect prompt injection show how untrusted context can override source boundaries and cause tool misuse or data leakage \cite{greshake2023not,yi2023bipia,zhan2024injecagent,ruan2024toolemu,debenedetti2024agentdojo}. In robotics, the analogous failure can become an unauthorized skill call, unsafe tool argument, or middleware command, but most evidence remains adjacent rather than directly VLA-validated. Toolformer, ReAct, Gorilla, ToolLLM, WebShop, and WebArena are useful here only as software-agent substrates; they should not be counted as VLA evidence unless robot skills or VLA-controlled actions are in the loop \cite{schick2023toolformer,yao2023react,patil2023gorilla,qin2023toolllm,yao2022webshop,zhou2023webarena}. ROS/ROS2 and robot cybersecurity work show that middleware, topics, services, node identity, and controller interfaces can be compromised \cite{quigley2009ros,dieber2020penetration,deng2022insecurityros2,quarta2017industrial,bonaci2015surgical,mayoral2020robot,yaacoub2022robotics,neupane2024security}. HRI work shows that robots can affect human trust, authority, and physical access decisions \cite{booth2017piggybacking,postnikoff2018robot,bainbridge2011benefits,saunderson2021persuasive}.

\textbf{Analogical evidence.} RAG and memory poisoning show how retrieved context can be corrupted and later reused \cite{zhong2023poisoning,deng2024pandora,zou2025poisonedrag,cho2024garag}. For VLA systems, the unvalidated but important hypothesis is multi-session embodied memory compromise: a poisoned object memory, map, user preference, or tool log could redirect future robot behavior. This remains mostly an open VLA attack problem, not a demonstrated field-wide result.
